%% ── Ⅱ. MATERIALS AND METHODS ───────────────────────────
\chapter{Main Body}

The main body generally consists of one to three chapters and describes the research contents, research methods, and research results. It should logically and clearly explain the process of answering the problem (research question) defined in the introduction.

\section{Research Methods}

\textbf{Function of the Research Methods}
\begin{itemize}
    \item To inform readers in detail about what the researcher conducted in order to answer the research questions described in the introduction
\end{itemize}

\textbf{Contents of the Research Methods}
\begin{enumerate}
    \item Materials (research subjects)
    \item Methods: detailed description of the methods used in the study, such as analytic, experimental, numerical, and statistical methods
\end{enumerate}

\textbf{Verb Tense in the Research Methods}
\begin{enumerate}
    \item Past tense should be used when describing performed methods (\ldots was measured using the OO method)
    \item Present tense should be used when describing data (\ldots is summarized as the mean and standard deviation)
\end{enumerate}

\section{Research Results}

\textbf{Function of the Research Results}
\begin{enumerate}
    \item To describe in detail the results obtained through the methods presented in the research methods section
    \item To support the research results with evidence such as data, tables, or figures
\end{enumerate}

\textbf{Contents Included in the Research Results}
\begin{enumerate}
    \item Only results relevant to the research questions raised in the introduction should be presented (data only without interpretation -- interpretation should be included in Chapter IV. Discussion)
    \item All results should be included regardless of whether they support the research hypothesis
\end{enumerate}

\textbf{Verb Tense in the Research Results}
\begin{itemize}
    \item Results should generally be reported in the past tense
\end{itemize}